% =====================================================================
%  THÈME 4 — Réactif limitant et avancement
%  Niveau DIFFICILE — Exercices
% =====================================================================
\documentclass[11pt,a4paper]{article}
\usepackage[utf8]{inputenc}
\usepackage[T1]{fontenc}
\usepackage{lmodern}
\usepackage[french]{babel}
\usepackage{amsmath,amssymb,mathtools}
\usepackage{icomma}
\usepackage[version=4]{mhchem}
\usepackage{siunitx}
\sisetup{output-decimal-marker={,},per-mode=symbol}
\DeclareSIUnit\Molar{\mole\per\litre}
\usepackage{xcolor}
\usepackage{tikz}
\usepackage{pgfplots}
\pgfplotsset{compat=1.18}
\usepackage[most]{tcolorbox}
\usepackage{enumitem}
\usepackage{array}
\usepackage{fancyhdr}
\usepackage[a4paper,margin=2.1cm,headheight=26pt,headsep=8pt]{geometry}

\definecolor{primary}{RGB}{142,93,168}
\definecolor{primarydark}{RGB}{82,45,108}
\definecolor{difcol}{RGB}{176,40,60}
\definecolor{reacA}{RGB}{0,90,200}
\definecolor{reacB}{RGB}{210,30,40}
\definecolor{prodC}{RGB}{0,140,90}
\newcommand{\nq}[1]{n_{\text{#1}}}
\newcommand{\Vm}{V_{\mathrm m}}

\pagestyle{fancy}\fancyhf{}
\fancyhead[L]{\small\textcolor{primarydark}{\textbf{Fiche 5} — Thème 4 : Limitant}}
\fancyhead[R]{\small\textcolor{difcol}{\textsc{Difficile}}}
\fancyfoot[C]{\small\thepage}
\renewcommand{\headrulewidth}{0.8pt}
\renewcommand{\headrule}{\hbox to\headwidth{\color{primary}\leaders\hrule height \headrulewidth\hfill}}

\newtcolorbox[auto counter]{exo}[1][]{enhanced,breakable,boxrule=1pt,arc=3pt,
  left=3mm,right=3mm,top=2.4mm,bottom=2.4mm,colback=difcol!5,colframe=difcol,
  fonttitle=\bfseries,coltitle=white,
  title={Exercice~\thetcbcounter\ \ifx\relax#1\relax\relax\else\textmd{--- \itshape #1}\fi}}

\setlength{\parindent}{0pt}
\setlength{\parskip}{3pt}

\begin{document}

\begin{center}
{\LARGE\bfseries\color{primarydark}Thème 4 — Réactif limitant}\\[1pt]
{\large\color{difcol}\textbf{Niveau Difficile}}\\[2pt]
{\itshape Objectif : lire un graphe d'avancement, traiter le cas stœchiométrique et les données mêlées.}\\
{\color{primary}\rule{0.64\linewidth}{1pt}}
\end{center}

\textbf{Données :} $\Vm=\SI{22,4}{\litre\per\mole}$ (TPN).

\begin{exo}[lecture d'un graphe d'avancement]
On étudie la réaction $\ce{2 A + B -> 3 C}$. Le graphe ci-dessous donne l'évolution des quantités de
matière en fonction de l'avancement $\xi$.
\begin{center}
\begin{tikzpicture}
\begin{axis}[
   width=10.5cm,height=5.8cm,
   xlabel={\small avancement $\xi$ (mol)},ylabel={\small $n$ (mol)},
   xmin=0,xmax=0.12,ymin=0,ymax=0.33,
   axis lines=left,
   xtick={0,0.02,0.04,0.06,0.08,0.10},
   ytick={0,0.05,0.10,0.15,0.20,0.25,0.30},
   grid=both,grid style={black!12},
   tick label style={font=\footnotesize},label style={font=\footnotesize},
   legend style={font=\scriptsize,at={(1.02,1)},anchor=north west,draw=black!25},
   clip=false,
]
  \addplot[reacA,line width=1.1pt] coordinates {(0,0.20)(0.10,0)}; \addlegendentry{$\nq{A}$}
  \addplot[reacB,line width=1.1pt] coordinates {(0,0.15)(0.10,0.05)}; \addlegendentry{$\nq{B}$}
  \addplot[prodC,line width=1.1pt] coordinates {(0,0)(0.10,0.30)}; \addlegendentry{$\nq{C}$}
\end{axis}
\end{tikzpicture}
\end{center}
\begin{enumerate}[label=\textbf{\alph*)},leftmargin=2em,topsep=1pt,itemsep=1.5pt]
  \item Lis sur le graphe les quantités initiales $\nq{A,0}$ et $\nq{B,0}$.
  \item Quel réactif s'annule le premier ? En déduire $\xi_{\max}$ et le réactif limitant.
  \item Donne les quantités finales de \ce{A}, \ce{B} et \ce{C} à $\xi_{\max}$.
  \item Parmi les trois courbes, laquelle a la pente la plus raide ? Relie ta réponse aux coefficients
        stœchiométriques.
\end{enumerate}
\end{exo}

\begin{exo}[le cas particulier stœchiométrique]
La formation de l'hydrure d'aluminium s'écrit $\ce{2 Al + 3 H2 -> 2 AlH3}$. On introduit
\SI{0,36}{\gram} d'aluminium et \SI{0,04}{\gram} de dihydrogène.
\emph{Masses molaires (\si{\gram\per\mole}) :} \ce{Al}=27 ; \ce{H2}=2 ; \ce{AlH3}=30.
\begin{enumerate}[label=\textbf{\alph*)},leftmargin=2em,topsep=1pt,itemsep=1.5pt]
  \item Convertis les deux masses en quantités de matière.
  \item Calcule les rapports $n/\nu$ pour \ce{Al} et \ce{H2}. Que constates-tu ?
  \item Un étudiant affirme : « le dihydrogène est le réactif limitant ». A-t-il raison ? Justifie.
  \item Quelle masse d'hydrure \ce{AlH3} obtient-on (rendement \SI{100}{\percent}) ?
\end{enumerate}
\end{exo}

\begin{exo}[données mêlées : masse, mole et solution]
Soit la réaction pondérée
\[ \ce{3 Co + 2 KNO3 + 8 HCl -> 2 NO + 3 CoCl2 + 2 KCl + 4 H2O}. \]
On introduit \SI{58,93}{\gram} de cobalt ($M(\ce{Co})=\SI{58,93}{\gram\per\mole}$), \SI{1,500}{\mole}
de \ce{KNO3} et \SI{20}{\milli\litre} d'acide chlorhydrique à \SI{6}{\Molar}.
\begin{enumerate}[label=\textbf{\alph*)},leftmargin=2em,topsep=1pt,itemsep=1.5pt]
  \item Convertis \emph{chaque} réactif en quantité de matière (attention aux unités !).
  \item Calcule les trois rapports $n/\nu$.
  \item Quel est le réactif limitant ? (Ce résultat servira de base au bilan complet du Thème 5.)
\end{enumerate}
\end{exo}

\end{document}
